\documentclass[11pt]{article}
\usepackage[utf8]{inputenc}
\usepackage[T1]{fontenc}
\usepackage{geometry}
\usepackage{booktabs}
\usepackage{longtable}
\usepackage{array}
\usepackage{tabularx}
\usepackage{float}
\usepackage{hyperref}
\usepackage{enumitem}
\usepackage{graphicx}
\geometry{margin=1in}
\hypersetup{colorlinks=true, linkcolor=blue, urlcolor=blue}
\newcolumntype{Y}{>{\raggedright\arraybackslash}X}

\title{Sound of Culture\\Phase 1.1 Report: Country-Only and Adjacent-Only Artist Universes}
\author{Detailed project report}
\date{\today}

\begin{document}
\maketitle
\tableofcontents
\clearpage

\section{Purpose of this report}
This report documents the retained workflow used to construct the artist-side
foundation of the project. Its job is not to describe the final chord-level
dataset. Instead, it explains how the project defined the \texttt{country-only}
artist frame, how it expanded that frame into \texttt{adjacent-only}, how
missing artist metadata were recovered, and which outputs are the canonical
deliverables for downstream work.

The key conceptual point is that Phase 1.1 solves an \emph{artist-frame}
problem before the project solves a \emph{song-frame} problem. The project
first needs a defensible answer to ``which artists belong to the core country
sample, which artists belong to the adjacent perimeter, and how much metadata
can be recovered for both groups?'' Only after that artist-side construction is
stable does the project move to Ultimate Guitar discovery and song-level
assembly.

\section{Final deliverables and how to read them}
The retained artist-side deliverables are:
\begin{itemize}[leftmargin=2em]
\item \texttt{artist\_universe\_country\_only.csv/.dta}: the canonical country
core used as the artist input for the Ultimate Guitar pipeline;
\item \texttt{artist\_universe\_adjacent\_only.csv/.dta}: the adjacent pool kept
separate from the core country frame;
\item \texttt{artist\_universe\_country\_plus\_adjacent.csv/.dta}: the final
deduplicated union of the two universes;
\item \texttt{artist\_universe\_qc\_report.md},
\texttt{artist\_universe\_methodological\_note.md}, and the coverage-estimate
files, which document quality, scope, and interpretation.
\end{itemize}

\begin{table}[H]
\centering
\caption{Canonical Phase 1.1 counts used in the retained workflow}
\begin{tabularx}{\textwidth}{Y r}
\toprule
Object & Count \\
\midrule
Base country-only artist universe & 2,333 \\
Adjacent-only artist universe & 2,493 \\
Country plus adjacent universe & 4,824 \\
Operational \texttt{artist\_universe\_country\_only.csv} after Billboard additions & 2,539 \\
\bottomrule
\end{tabularx}
\end{table}

The distinction between the first and the fourth row matters. The retained
Phase 1.1 universe construction produces the clean artist-side core
(\texttt{2,333}). Later, the Phase 1.2 Billboard supplement adds artists needed
to recover missing chart material, which is why the current operational
\texttt{artist\_universe\_country\_only.csv} is larger (\texttt{2,539}). This
report therefore treats the clean country-only universe as the canonical
construction object and the enlarged country-only file as the operational input
handed over to Phase 1.2.

\section{Retained workflow}

\subsection{Stage 1. Build the base country-artist datasets}
The first retained construction layer is the base country-artist build in
\texttt{code/build\_country\_artists\_dataset.py}. The key orchestration lives in:
\begin{itemize}[leftmargin=2em]
\item \texttt{prepare\_dataset()} at lines 2123--2258;
\item \texttt{deduplicate()} at lines 2329--2353;
\item \texttt{write\_qc\_report()} at lines 2393--2439;
\item \texttt{main()} at lines 2442--2465 and the code that follows in the same
entrypoint.
\end{itemize}

This builder does five things in a disciplined order.
\begin{enumerate}[leftmargin=2em]
\item It collects seed candidates from the Country Music Hall of Fame, from
Wikipedia pages and categories, and from broad country-related Wikidata
queries.
\item It merges the seed layer with structured Wikidata details.
\item It enriches identity and geography from Wikipedia when the structured
layer is incomplete.
\item It applies targeted MusicBrainz enrichment when important key fields
remain unresolved.
\item It deduplicates, assigns sample flags, and exports the country-artist
tables in both \texttt{csv} and \texttt{dta}.
\end{enumerate}

The important methodological choice here is the source hierarchy. The retained
workflow does not begin with weak web scraping. It begins with structured
artist identifiers and only later uses semi-structured or targeted fallback
sources.

\subsection{Stage 2. Build the country core and the adjacent-only perimeter}
The second retained layer is the integrated artist-universe build in
\texttt{code/build\_extended\_artist\_universe.py}. The main retained blocks are:
\begin{itemize}[leftmargin=2em]
\item \texttt{build\_core\_universe()} at lines 571--583;
\item \texttt{build\_adjacent\_seed\_pool()} at lines 586--607;
\item \texttt{resolve\_adjacent\_seeds()} at lines 610--691;
\item \texttt{build\_adjacent\_universe()} at lines 694--845;
\item \texttt{write\_qc\_report()} at lines 869--892;
\item \texttt{write\_method\_note()} at lines 895--945;
\item \texttt{main()} at lines 948 onward.
\end{itemize}

This script is where the final artist-universe logic becomes clear. The country
core is not rebuilt from scratch. Instead, it is integrated from three existing
project assets:
\begin{itemize}[leftmargin=2em]
\item artists already observed in the current song-side file;
\item artists already present in \texttt{country\_artists\_restricted};
\item artists already present in \texttt{country\_artists\_master}.
\end{itemize}

After the core is stabilized, the script builds an adjacent seed pool from
genre-oriented sources and then resolves those candidates through the same
artist-enrichment logic used on the country side. The adjacent perimeter is not
treated as ``country by another name.'' It is retained as a separate staging
universe that widens later discovery potential without collapsing genre
identity into a single bucket.

The adjacent seed logic is materially important. The retained adjacent tags are
\texttt{folk}, \texttt{americana}, \texttt{bluegrass},
\texttt{rockabilly}, \texttt{western swing}, \texttt{outlaw country},
\texttt{roots rock}, and \texttt{country gospel}. The script then classifies
adjacent cases into buckets and keeps the \texttt{adjacent-only} pool separate
from the core country frame.

\subsection{Stage 3. Recover missing metadata on the country-only file used downstream}
The third retained layer is a direct top-up of the exact
\texttt{artist\_universe\_country\_only.csv} file that Phase 1.2 will use. The
retained logic is in
\texttt{code/enrich\_artist\_universe\_missing\_metadata.py}, especially
\texttt{enrich\_country\_only\_direct()} at lines 2062--2085 and
\texttt{main()} at lines 2119 onward.

This pass exists because the project separates two tasks:
\begin{itemize}[leftmargin=2em]
\item building a broad artist-side frame;
\item making sure the exact country-only file used by the Ultimate Guitar
pipeline is as complete as possible on the key demographic variables.
\end{itemize}

The retained enrichment order is explicit and hierarchical:
\begin{enumerate}[leftmargin=2em]
\item parse and clean the geography already present;
\item resolve missing Wikidata QIDs from names;
\item enrich from targeted Wikidata lookups;
\item enrich from Wikipedia sources;
\item resolve missing MusicBrainz MBIDs;
\item run targeted MusicBrainz enrichment;
\item run generic MusicBrainz enrichment;
\item recover origin-year information from Wikipedia for bands and residual
cases;
\item use Discogs and iTunes as additional fallback sources;
\item use official websites when structured public evidence exists;
\item apply web-confirmed manual overrides;
\item clean geography again and infer state from already-observed
city-country combinations;
\item refresh derived columns and save the final country-only file.
\end{enumerate}

\section{Source hierarchy and variable recovery}
The retained workflow deliberately uses a source hierarchy rather than a flat
merge of sources.

\subsection{Primary source layer}
The first layer is structured and semi-structured artist data:
\begin{itemize}[leftmargin=2em]
\item Wikidata for identifiers, birth dates, places, genres, occupations, and
citizenship;
\item Wikipedia infoboxes and lead text for recovery when structured fields are
missing or incomplete;
\item MusicBrainz for artist identity, area, and lifespan-style metadata.
\end{itemize}

\subsection{Secondary fallback layer}
Only after that first layer does the project use:
\begin{itemize}[leftmargin=2em]
\item Discogs;
\item iTunes / Apple Music;
\item official websites via structured page metadata;
\item targeted manual overrides when public evidence is clear enough.
\end{itemize}

\subsection{Priority variables}
The project pushed hardest on:
\begin{itemize}[leftmargin=2em]
\item \texttt{name\_primary};
\item \texttt{birth\_year};
\item \texttt{birth\_city};
\item \texttt{birth\_state};
\item \texttt{birth\_country}.
\end{itemize}

These variables matter for the later substantive design, because they are the
artist-side bridge to time, geography, and eventual state-level cultural
measurement.

\begin{table}[H]
\centering
\caption{Residual missingness in the current operational country-only file}
\begin{tabularx}{\textwidth}{Y r}
\toprule
Variable & Remaining missing \\
\midrule
\texttt{birth\_year} & 118 \\
\texttt{birth\_city} & 178 \\
\texttt{birth\_state} & 122 \\
\texttt{birth\_country} & 96 \\
\texttt{wikidata\_qid} & 89 \\
\texttt{musicbrainz\_mbid} & 531 \\
\bottomrule
\end{tabularx}
\end{table}

The retained principle is that the project prefers reliable missingness over
fabricated certainty. If a field cannot be supported by sufficiently credible
evidence, it remains missing.

\section{Final outputs kept for downstream work}
The artist-side outputs that are actually retained are:
\begin{itemize}[leftmargin=2em]
\item the three final universe files in \texttt{csv} and \texttt{dta};
\item the quality-control report;
\item the methodological note on denominators and coverage;
\item the coverage-estimate files;
\item the replication package for exact restoration of the final universe
deliverables.
\end{itemize}

This last point matters. The retained project workflow is not just ``run the
live builder again.'' It also preserves a stable replication package. The exact
artist-universe restoration wrapper is in
\texttt{code/run\_artist\_universe\_replication.py}, with the package structure
and restore/copy logic at lines 6--20, 57--65, 68--75, 87--108, 111--121, and
128--145.

\section{Downstream artist-origin repair in the final song dataset}
Phase 1.1 also has a downstream implication that only became visible once the
Phase 2 exploratory plots were produced. The final song-level dataset inherits
artist geography from the country-only universe and from later merge layers, so
the project needed a second artist-side repair pass directly on the final
dataset once the exploratory figures revealed too many residual
\texttt{Unknown} macro regions.

The retained downstream repair now lives in
\path{execution/phase_01_dataset_construction/backfill_country_only_final_artist_metadata.py},
especially the suspicious-match reset and manual override layer at lines
107--265, the conservative lookup logic at lines 379--523, and the save path at
lines 526--579.

The important retained principle is conservative correction rather than blind
completion. The script first clears inherited false positives for ambiguous
artist strings, then reapplies only stricter matches, and finally uses a small
set of web-confirmed overrides for the cases where public evidence is clear
enough. This is why some rows that previously looked ``complete'' were allowed
to become missing again: the project now prefers verified missingness to
incorrect country or state assignments.

\begin{table}[H]
\centering
\caption{Current artist-origin status inside the final song dataset after the retained repair pass}
\begin{tabularx}{\textwidth}{Y r}
\toprule
Metric & Count \\
\midrule
Rows still missing \texttt{birth\_state} & 133 \\
Rows still coded as missing or \texttt{Unknown} in \texttt{us\_macro\_region} & 122 \\
Rows explicitly coded as \texttt{Non-US} in \texttt{us\_macro\_region} & 2,768 \\
Rows still missing \texttt{birth\_country} & 122 \\
Unique artists still missing \texttt{birth\_state} & 39 \\
Unique artists still missing or \texttt{Unknown} in \texttt{us\_macro\_region} & 35 \\
Unique artists explicitly coded as \texttt{Non-US} in \texttt{us\_macro\_region} & 96 \\
Unique artists still missing \texttt{birth\_country} & 35 \\
US artists still missing \texttt{birth\_state} & 0 \\
\bottomrule
\end{tabularx}
\end{table}

The last row matters most for the substantive geography problem. After the
retained repair pass and the additional curated update on 2026-04-20, no
artist that is currently coded as US-origin still lacks a state assignment in
the final song-level file. \texttt{The Wreckers} now receives a curated
group-origin assignment of \texttt{Nashville, Tennessee}, while non-US artists
are explicitly coded as \texttt{Non-US} rather than being left under
\texttt{Unknown}. The remaining \texttt{Unknown} block is therefore much
smaller and now tracks genuinely unresolved geography.

\section{Retained versus non-retained choices}
\subsection{Retained choices}
The following choices are retained in the final workflow:
\begin{itemize}[leftmargin=2em]
\item keeping \texttt{country-only} and \texttt{adjacent-only} as two separate
objects rather than collapsing them into a single artist list;
\item using structured sources first, then semi-structured fallbacks, then
targeted web-confirmed overrides only when necessary;
\item treating the country-only file as the operational input for song-side
discovery;
\item preserving exact-replication packaging alongside the live build logic.
\end{itemize}

\subsection{Not retained as final choices}
The following ideas were not retained as the final workflow:
\begin{itemize}[leftmargin=2em]
\item using the full adjacent universe directly as the scrape input for the
main Ultimate Guitar country dataset;
\item forcing complete metadata fill when public evidence remained weak;
\item redefining adjacent artists as if they were equivalent to the core
country frame.
\end{itemize}

Those choices were rejected because they would either blur the substantive
meaning of the country sample or create unreliable metadata.

\section{Exact scripts used in the retained workflow}
\scriptsize
\begin{longtable}{p{0.34\textwidth}p{0.18\textwidth}p{0.35\textwidth}}
\toprule
Script copy in this folder & Main line ranges used & Role in the retained workflow \\
\midrule
\endfirsthead
\toprule
Script copy in this folder & Main line ranges used & Role in the retained workflow \\
\midrule
\endhead
\path{code/build_country_artists_dataset.py} & 2123--2258; 2329--2353; 2393--2439; 2442 onward & Build the base country-artist datasets, deduplicate, and export the first QC layer. \\
\path{code/build_extended_artist_universe.py} & 571--583; 586--607; 610--691; 694--845; 869--945; 948 onward & Build the country core, expand to adjacent-only, and write the final universe QC and methodological notes. \\
\path{code/enrich_artist_universe_missing_metadata.py} & 2062--2085; 2119 onward & Recover missing key metadata directly on the operational country-only file used downstream. \\
\path{code/run_artist_universe_replication.py} & 6--20; 57--65; 68--75; 87--108; 111--121; 128--145 & Restore or refresh the exact artist-universe replication package. \\
\bottomrule
\end{longtable}
\normalsize

\section{Stata note}
There is no retained substantive Stata workflow in Phase 1.1. The artist-side
construction is Python-based. Stata launchers exist only as thin replication
wrappers, so this report does not include Stata code snippets.

\section{Conclusion}
Phase 1.1 gives the project a disciplined artist-side frame before any
song-level scraping begins. The core contribution of this phase is not only the
three final universe files. It is the combination of:
\begin{itemize}[leftmargin=2em]
\item a separated country and adjacent design;
\item a documented source hierarchy for metadata recovery;
\item a clear rule that country-only is the operational scrape input;
\item a reproducible package that restores the final artist-universe outputs.
\end{itemize}

That is the retained artist-side workflow that the rest of the project now
inherits.

\end{document}
